%!TEX TS-program = xelatex
\documentclass[]{friggeri-cv}
\usepackage{afterpage}
\usepackage{hyperref}
\usepackage{csquotes}
\usepackage[spanish]{babel}
\usepackage{color}
\usepackage{xcolor}
\hypersetup{%
    pdftitle={Curriculum Vitae},
    pdfauthor={Arturo Volpe},
    pdfsubject={Resumen de vida Profesional},
    pdfkeywords={},
    colorlinks=false,       % no link border color
   allbordercolors=white    % white border color for all
}
\addbibresource{bibliography.bib}
\RequirePackage{xcolor}
\definecolor{pblue}{HTML}{0395DE}

\begin{document}
\header{Arturo }{Volpe}
{Ingeniero en Informática}

% Fake text to add separator
\fcolorbox{gray}{gray}{\parbox{\dimexpr\textwidth-40\fboxsep-2\fboxrule}{%
.....
}}

% In the aside, each new line forces a line break
\begin{aside}
	\section{}~
    \section{}~
    \section{}~
    \section{Celular}
        +595972105788
        ~
    \section{Correo}
        \href{mailto:arturo@volpe.com.py}{arturo@volpe.com.py}
        \href{mailto:arturovolpe@gmail.com}{arturovolpe@gmail.com}
        ~
    \section{Repositorios}
        \href{https://github.com/avolpe}{github.com/avolpe}
        \href{https://gitlab.com/avolpe}{gitlab.com/avolpe}
        ~
    \section{Blog}
        \href{https://www.volpe.com.py}{volpe.com.py}
        ~
\section{Idiomas}
\textbf{Español}\includegraphics[scale=0.40]{img/5stars.png}
\textbf{Ingles}\includegraphics[scale=0.40]{img/3stars.png}
\end{aside}


\section{Experiencia}
\begin{entrylist}
\entry
    {2017 - 2018}
    {Codium S.A.}
    {}
    { \proyecto
        {Bancard S.A.}
        {2017-2018}
        {}
        {Consultor}
        {Soporte para implementación de un esquema de integración continua para los proyectos satélite
    de Bancard, lo que incluye sistemas desarrollados sobre Ruby on Rails, y Spring.
    
    Se utilizo Jenkins para el manejo de tareas, geminabox y artifactory para la gestión de librerías,
    SonarQube para análisis estático y Gitlab como repositorio de código fuente.} 
        {170}
    }
    {}
\entry
    {2015-2017}
    {\href{http://www.cds.com.py}{Centro de Desarrollo Sostenible S.A.}}
    {}
    {\proyecto
        {Financiera el Comercio}
        {2016-2017}
        {}
        {Arquitecto de Software}
        {\textbf{Arquitectura y mantenimiento} de sistemas satélite de la Financiera, incluyendo BancaWeb, Aplicaciones móviles, sistemas de gestión internos e integración con sistemas Externos. Encargado de proceso de 
        \textbf{Integración continua} a través de la automatización del proceso de construcción de todos los sistemas y automatización de pruebas de integración y e2e. En el área de \textbf{producción}, automatización del proceso de monitoreo de sistemas y generación de alertas.
        } 
        {976}
    
    \proyecto
        {\href{https://www.mytappweb.com}{myTapp}}
        {2016-2017}
        {Plataforma para gestión de comercios y pedidos con integración a plataformas de pago.}
        {Arquitecto y full stack}
        {
        Desarrollo de aplicación móvil desarrollada con  \textbf{ionic (v1)}, 
        backend utilizando el stack JavaEE~(jax-rs, JPA, EJB) y una base de datos
        PostgreSQL.
        
		Además implemente un esquema de \textit{continuous delivery} para este proyecto,
        utilizando \textit{Jenkins, Rancher, Docker, Arquillian, SonarQUBE, AWS y Scaleway}. 
        
        En producción se utilizaron  otros sistemas complementarios de código abierto 
        como son: \textit{PictShare, pictfit} para la gestión de imágenes, redis para una
        cache de alto rendimiento y nginx como proxy reverso.
        
        El sistema en producción funciona sobre AWS (RDS, y instancias EC2 manejadas a 
        través de un Rancher). Además incluye un sistema de monitoreo del sistema
        utilizando \textit{Graphite y Dropwizard}.
        } 
        {}
        
    \proyecto
        {\href{http://www.cds.com.py/productos/kinematics/}{Kinematic}}
        {2016-2017}
        {Plataforma automatizada de entrenamiento de tiro de balón}
        {Arquitecto y full stack}
        {Consta de dos componentes principales, un \textbf{core} de detección de movimiento de pelotas en movimiento desarrollado en python y c++ utilizando OpenCV. Así mismo, un componente de \textbf{front-end}, el cual se encuentra desarrollado con python utilizando django.}
        {300}
        
    \proyecto
        {\href{https://mawio.net/}{MAWIO}}
        {2016-2017}
        {Mapa Web Interactivo de Obras de la Essap}
        {Arquitecto y full stack}
        {SPA desarrollada utilizando principalmente en el front-end \textbf{angular, leaflet y jQuery datatables}, en el back-end \textbf{Jaxa EE (jax-rs, JPA y EJB),}, una base de datos \textbf{PostgreSQL con PostGIS}, y manejo del entorno de desarrollo y producción utilizando \textbf{docker, nginx, wildfly 9 y redis}}
        {300}
        
        
      }
      {}
    
  \entry
    {2015-2017}
    {\href{http://www.cds.com.py}{Centro de Desarrollo Sostenible S.A.}}
    {}
    {
    \proyecto
        {\href{http://www.cds.com.py/productos/sistemas-de-gestion/sistema-de-gestion-hospitalaria/}{Sistema de Gestión Hospitalaria}}
        {2015-2017}
        {Sistema integral para hospitales}
        {Arquitecto y líder de proyecto}
        {SPA desarrollada utilizando principalmente en el front-end 
        \textbf{angular, boostrap-ui y jQuery datatables}, en el back-end 
        \textbf{Jaxa EE (jax-rs, JPA y EJB)}, una base de datos PostgreSQL, 
        y manejo del entorno de desarrollo y producción utilizando 
        \textbf{docker, nginx, wildfly 9 y redis}}
        {1920}
        
    \proyecto
        {Portal de datos abiertos para el Poder Judicial}
        {2015}
        {Portal de datos abiertos que incluye varias fuentes de datos.}
        {full stack}
        {Portal de datos abiertos desarrollado como una SPA utilizando 
        \textbf{angular, boostrap-ui y jQuery datatables}, en el back end 
        se utilizo el stack estándar de Java EE (jax-rs, JPA, EJB), como base 
        de datos se utilizo Microsoft SQL Server 2008.}
        {352} 
        
    \proyecto
        {\href{http://bacheando.com/}{Bacheando}}
        {2015}
        {Proyecto de localización de incidentes urbanos}
        {full stack}
        {Permite la geolocalización de baches, perdidas, hundimientos y
            cortes de la municipalidad de Asunción.
            Posee un componente web (desarrollado utilizando 
            \textbf{angular, d3, boostrap y leaflet}), así como una 
            aplicación móvil para dispositivos Android y un servidor desarrollado 
            utilizando wildfly y Java EE(jax-rs, JPA).}
        {176} 
        
	}
    {}


\entry
    {2012 - 2015}
    {Facultad Politécnica}
    {Arquitecto de Software}
    {Ingrese como programador junior, desempeñando cargos como la creación de
        aplicaciones base para el desarrollo, tareas de diseño de interfaces
        como de lógica de negocio, integración de componentes, e investigación
        de tecnológicas.

        Proyectos:
        \begin{itemize}
            \item Sistema Integral de Gestión Hospitalaria (Hospital de
                Clínicas)
            \item Sistema de Gestión Convergente de Nueva Generación (Copaco)
        \end{itemize}}
    {5280}
\entry
    {2011 - 2012}
    {Konecta S.A.}
    {Programador Junior}
    {Como programador junior, mis responsabilidades incluían el diseño y
        desarrollo de software para dispositivos móviles, y la creación de
        servicios webs, tanto servicios que interactuaban directamente con
        aplicaciones móviles, como otros que interactuaban con portales webs
        \\}
    {1408}
\entry
    {2011 - 2011}
    {Global S.A.}
    {Programador Junior}
    {Como programador junior mis funciones iban desde correcciones y pruebas de
        aplicaciones en fase final de desarrollo, hasta realizar funciones
        relacionadas con el mantenimiento de la base de datos y su correcto
        funcionamiento. Además de administración de servidores. \\}
    {300}
\entry
    {2010 - 2011}
    {Bunker Games}
    {Programador}
    {Como programador principal y único, mis responsabilidades iban desde la
        correcto funcionamiento del videojuego, la integración del entorno
        gráfico, programación de la inteligencia artificial, cálculo de
        recorridos, optimización, etc. \\}
    {1408}

\end{entrylist}

\section{Freelancer}
\proyectof
    {2018}
    {\href{https://www.pemoi.com.py}{Pemoi}}
    {\href{https://www.pemoi.com.py/}{Yo Dono}}
    {Plataforma de Crowfounding para campañas políticas.}
    {Consiste en una aplicación multiplataforma~(Andorid e iOS) desarrollada utilizando
     Ionic 3, servicios de firebase (social login, realtime database, deep links, functions, 
     storage, etc), fabric.io (beta, crashlytics, fastlane, etc).
     
     Incluye un sistema de pagos digitales utilizando la plataformas de
     Tigo Money, Bancard S.A. y \href{https://www.pagopar.com.py}{Pagopar}.
     
     }
    {300}
\proyectof
    {2015}
    {\href{https://apps.facebook.com/pechugon_pepe}{PepeChugon Game}}
    {\href{http://www.animata.com.py/}{Animata}}
    {Consiste en un juego multiplataforma~(Android, Facebook Canvas y Windows)
        y un servidor donde se almacena la información y el rendimiento de los
        usuarios. El juego fue desarrollado para la campaña publicitaria del día
        del niño.}
    {El Juego fue desarrollado utilizando Unity3D y un componente de backend 
    desarrollado utilizando OpenShift con PostgreSQL.}
    {300}
        

\section{Experiencia académica}
\clase
    {2016}
    {Curso de Verano}
    {\href{http://www.pol.una.py/cursosverano/index.php?option=com_content&view=article&layout=edit&id=91}{Innovación
            social y transparencia mediante datos abiertos}} 
    {Auxiliar}
    {FP-UNA}
    {El curso busca crear una comunidad de personas que puedan aportar ideas
        innovativas con alto impacto social y que fomenten la transparencia en
        el uso de los recursos públicos. }
\clase
    {2016}
    {Licenciatura en ciencias Informáticas}
    {Algoritmos I y II} 
    {Auxiliar de enseñanza}
    {FP-UNA}
    {Introducción a algoritmos, conceptos básicos y programación con SLE2 y C.}

\clase
    {2016}
    {Ingeniería en informática}
    {Gestión de Centro de Cómputos} 
    {Auxiliar de enseñanza}
    {FP-UNA}
    {}

\section{Educación}
\begin{entrylist}
  \educacion
    {2008 - 2015}
    {Ingeniero en Informática}
    {Facultad Politécnica - Universidad Nacional de Asunción}
    {\textbf{Tesis de grado:} Juegos serios como apoyo a la formación de profesionales: Una aplicación al área de enfermería}
  \educacion
    {2005 - 2007}
    {Bachiller Técnico en Informática}
    {Colegio de San José}
    {}
\end{entrylist}

\begin{absolutelynopagebreak}
\section{Referencias}
\begin{entrylist}
    \educacion{M.Sc.}{Martin Abente}{}{}
    \educacion{M.Sc.}{Guillermo Gonzalez}{}{}
    \educacion{M.Sc.}{Maria Elena Garcia}{}{}
\end{entrylist}
\end{absolutelynopagebreak}

\begin{flushleft}
    \emph{\today{}}
\end{flushleft}
\begin{flushright}
    \emph{Arturo Volpe}
\end{flushright}



\end{document}
